%!TEX root = ../template.tex
%%%%%%%%%%%%%%%%%%%%%%%%%%%%%%%%%%%%%%%%%%%%%%%%%%%%%%%%%%%%%%%%%%%%
%% chapter5.tex
%% NOVA thesis document file
%%
%% Chapter with Validation and Testing
%%%%%%%%%%%%%%%%%%%%%%%%%%%%%%%%%%%%%%%%%%%%%%%%%%%%%%%%%%%%%%%%%%%%


\typeout{NT FILE chapter5.tex}

\chapter{Verification and Testing}
\label{cha:verificaton_testing}

In this chapter we'll go over the verification and testing procedures conducted to evaluate the proposed wearable monitoring system and the \gls{ML} models. This evaluation is divided into three sections addressing the full \gls{IoT} system.
Starting with the hardware, we examine the position of the prototypes, and how they are attached, and the anatomical placement of the \gls{EMG} electrodes on the Gastrocnemius Lateralis. Secondly, an in-depth statistical analysis of the trained 
\gls{ML} models is conducted. In this section, by leveraging classification metrics, confusion matrices, and feature importance distributions, both developed \gls{ML} models' performances are compared against each other and evaluated. Here, 
we also determine the efficacy of fusing \gls{IMU} kinematics with \gls{EMG} neuromuscular data for basketball movements classification. In the final section, the system's performance is reviewed by assessing both models' latency and the resilience 
to asynchronous \gls{BLE} transmissions.

\section{Data Collection}
\label{sec:data_collection}

The data sets used to train and test the models were collected from several volunteer athletes. They were asked to do each of the five movements - Idle, Walking, Running, Defensive Slide and Rebounding - for some minutes 
each. The cramping prevention system was only tested on one player that pushed his muscles to near involuntary contractions.


\subsection{Prototype Placement and Configuration}
\label{ssec:prototype_placement}

The prototypes were mounted directly in the athletes' footwear to approximate the ideal positioning (on the bottom of the feet as an insole or a sock). The devices were secured using the shoe laces and the powerbank was secured in the users 
socks, utilizing the existing infrastructure. With this configuration, the \gls{IMU} accurately captures the kinematics of the foot, such as accelerations and rotational velocities during the athletes' movements. Furthermore, this setup 
minimizes the mechanical strain exerted by the power cables on the USB ports, while keeping the \gls{EMG} leads in close proximity to the chosen muscle group.


\begin{figure}[htbp]
  \centering
  \includegraphics[width=0.5\linewidth]{deviceInPlace.jpg}
  \caption{An athlete wearing the prototype}
  \label{fig:protoInPlace}  
\end{figure}


The muscle chosen to be measured was the Gastrocnemius Lateralis as it is one of the most important muscles for explosive jumping motion.
The position for the \gls{EMG} sensors was chosen as recommended by the \gls{EU}'s SENIAM (Surface ElectroMyoGraphy for the Non-Invasive Assessment of Muscles) project (~\cite{emgPosition}) ensuring standardization and maximization of the 
signal-to-noise ratio.

\begin{figure}[htbp]
  \centering
  \includegraphics[width=0.6\linewidth]{emgPosition.jpg}
  \caption{Gastrocnemius Lateralis Positioning ~\cite{emgPosition}}
  \label{fig:emgPosition}
\end{figure}

%\subsection{Collecting Data and Testing}
%\label{ssec:prototype_placement}

\section{Machine Learning Model Evaluation}
\label{sec:ml_evaluation}

This section goes over an evaluation of the two \gls{ML} models trained for this study. To effectively assess their performance, stability, and learning behaviour, we will examine the graphics generated throughout 
the training phase giving us clear insights into how each model converges, identifying potential issues like under or overfitting, and ultimately compare their suitability for the objectives of this project.


\subsection{Training Metrics and Performance Curves}           %Rever
\label{ssec:training_metrics}

For the TinyML model, the learning process is monitored iteratively across the training epochs. Figure~\ref{fig:tinyMLAccuracy} illustrates the model's accuracy progression, comparing the training dataset against the 
validation dataset. This behaviour is further detailed by the loss curve in Figure~\ref{fig:tinyMLLoss}, which tracks the error rate during training. 




%To evaluate the learning stability and generalization capacity of the edge-tier TinyML neural network, the classification accuracy and categorical cross-entropy loss were continuously logged across the 40 training epochs. Evaluating these 
%temporal learning curves is critical to ensure the lightweight model converged properly without memorizing localized noise within the dataset.
%As illustrated in \autoref{fig:tinyMLAccuracy}, the network demonstrated rapid initial learning, with both training and validation accuracy climbing steeply during the first 15 epochs. The model successfully converged, achieving a final 
%validation accuracy of approximately 95\%. The absence of a severe divergence between the training and validation accuracy curves indicates that the model did not suffer from substantial overfitting, despite the high-dimensional 450-feature 
%input vectors. 
%This stability is further corroborated by the loss metrics depicted in \autoref{fig:tinyMLLoss}. The categorical cross-entropy loss for the validation split decayed smoothly alongside the training split, plateauing effectively after 20 epochs. 
%The minor fluctuations observed in the later epochs are characteristic of Adam optimizer adjustments across a complex, multi-class distribution, but the overall downward trajectory confirms that the TinyML model successfully captured the 
%underlying biomechanical and neuromuscular patterns required for on-device inference.
%A steady increase and subsequent plateau in validation accuracy suggest successful convergence.
%By analysing the gap between the training and validation loss, we can assess the model's generalization capabilities.


\begin{figure}[htbp]
  \centering
  \includegraphics[width=0.9\linewidth]{TinyMLAccuracy.jpg}
  \caption{TinyMLAccuracy}
  \label{fig:tinyMLAccuracy}
\end{figure}

\begin{figure}[htbp]
  \centering
  \includegraphics[width=0.9\linewidth]{TinyMLLoss.jpg}
  \caption{TinyMLLoss}
  \label{fig:tinyMLLoss}
\end{figure}


\begin{figure}[htbp]
  \centering
  \includegraphics[width=0.9\linewidth]{RFROC.jpg}
  \caption{RF ROC Curve}
  \label{fig:RFROC}
\end{figure}

\begin{figure}[htbp]
  \centering
  \includegraphics[width=0.9\linewidth]{RFPrecisionRecall.jpg}
  \caption{RF Precision Recall Curve}
  \label{fig:RFPrecision}
\end{figure}


\subsection{Feature Importance and Sensor Fusion}
\label{ssec:feature_importance}

The next Figure~\autoref{fig:featureImportance} shows that even the \gls{IMU} data has greater importance and the \gls{EMG} provides a very small contribution to the movement identification (+/-10\%). This means that 
for better performance the \gls{EMG} data could be cut and be used only for the injury prevention algorithm.

\begin{figure}[htbp]
  \centering
  \includegraphics[width=0.9\linewidth]{featureImportance.jpg}
  \caption{Feature Importance Graph}
  \label{fig:featureImportance}
\end{figure}

\subsection{Confusion Matrices and Error Analysis}
\label{ssec:confusion_matrices_ML}

\begin{figure}[htbp]
  \centering
  \includegraphics[width=0.9\linewidth]{tinyMLConfusionMatrix.jpg}
  \caption{TinyML Confusion Matrix}
  \label{fig:tinyMLConfusionMatrix}
\end{figure}


\begin{figure}[htbp]
  \centering
  \includegraphics[width=0.9\linewidth]{RFConfusionMatrix.jpg}
  \caption{Random ForestL Confusion Matrix}
  \label{fig:RFConfusionMatrix}
\end{figure}




\section{System Performance} %rever
\label{sec:system_performance}

A live test was performed in similar fashion to the data collection phase with the simple difference of running the server and taking note of the predictions and the movement made by the athlete (using the "Send Predictions" switch). The 
player was also pushed until the fatigue warning triggered.


\subsection{Overall System Accuracy}
\label{ssec:overall_accuracy}

The primary objective of this test was to compare the classification accuracy of the ultra-lightweight Edge model against the deeper Cloud model. In terms of latency, the tinyML model responded instantly while the \gls{RF} model had a slight
 delay.


\begin{figure}[htbp]
  \centering
  \includegraphics[width=0.8\linewidth]{testOverallAccuracyComparison.jpg}
  \caption{Overall Classification Accuracy: Local TinyML vs. Cloud Random Forest}
  \label{fig:overall_accuracy}
\end{figure}

As illustrated above (~\autoref{fig:overall_accuracy}), the localized TinyML edge inference successfully outperformed the cloud-based Random Forest model, achieving an overall accuracy of 74.65\% compared to the server's 71.11\%. Even tho 
the local model achieved better results, with a bigger dataset the server should outperform it, especially with more confusing movements.

\subsection{Confusion Matrices and Error Analysis}
\label{ssec:confusion_matrices_test}

In the graph below, we can see that overall, besides running, the TinyML model either outperformed the \gls{RF} model or was nearly as accurate. The \gls{RF} model struggled to correctly identify the movements 
with less data (Running, Rebounding and Defensive Sliding) but excelled at the Idle and Walking movements. 

\begin{figure}[htbp]
  \centering
  \includegraphics[width=0.9\linewidth]{testMovementAccuracy.jpg}
  \caption{Classification Accuracy by Movement}
  \label{fig:movement_accuracy}
\end{figure}

Continuing the system's assessment, the confusion matrices for both models were generated:
 
\begin{figure}[htbp]
  \centering
  \includegraphics[width=1\linewidth]{testConfusionMatrixRF.jpg}
  \caption{Confusion Matrix for the Random Forest Model}
  \label{fig:confusion_matrix_RF}
\end{figure}

In the \gls{RF}'s matrix we can see that the model tended to predict the movement Walking or Idle while barely choosing Rebound or Defensive Slide. This might be happening due to the low volume of data of these two 
movements when compared to the others.
The Rebound movement was very confusing for the model predicting it to be Walking half the time and barely getting it right.

\begin{figure}[htbp]
  \centering
  \includegraphics[width=1\linewidth]{testConfusionMatrixTinyML.jpg}
  \caption{Confusion Matrix for the TinyML Model}
  \label{fig:confusion_matrix_TinyML}
\end{figure}

On the other hand, the TinyML model mixed the Running movement with the Defensive Sliding and the Rebound movement 

These graphics reveal that both models and a near-perfect precision when dealing with static or low-intensity states (Idle and Walking) while the more explosive movements


%The matrices reveal that both models demonstrate near-perfect precision when identifying static or low-intensity states, such as ``Idle'' and ``Walking''. However, transient and explosive movements introduced a higher degree of misclassification. 
%For instance, the models occasionally confused a ``Defensive Slide'' with ``Running'', as both movements generate high-amplitude, lateral acceleration vectors on the \gls{IMU}. The TinyML model proved slightly more robust at isolating the 
%``Defensive Slide'' signature, contributing to its higher overall accuracy.

\subsection{Movements Performance Breakdown}
\label{ssec:movement_accuracy}

%Finally, analysing the accuracy on a per-movement basis highlights the specific biomechanical strengths of the neural network architecture. 



%As seen in \autoref{fig:movement_accuracy}, the Edge TinyML model demonstrated superior or equal performance across almost all dynamic classes. Notably, during the ``Rebound'' phase—a highly complex, multi-axial movement involving a sudden 
%vertical leap and harsh landing—the Edge model maintained a strong classification rate, whereas the Cloud model's accuracy dropped significantly. This suggests that the deep neural network layers of the TinyML model were better suited to 
%capturing the non-linear kinematic relationships of explosive jumping than the decision trees utilized by the Random Forest algorithm.


\section{Discussion of Results}
\label{sec:discussion}

